%-----------------------------------------------------------------------------
%  Bernoulli-EKF Multi-Target Tracker on SE(3): A Minimal PMBM Reduction
%-----------------------------------------------------------------------------
\documentclass[11pt]{article}
\usepackage[margin=1in]{geometry}
\usepackage{amsmath,amssymb,amsthm,mathtools,bm,enumitem,hyperref}

\DeclareMathOperator{\Tr}{tr}
\DeclareMathOperator{\Exp}{Exp}
\DeclareMathOperator{\Log}{Log}
\DeclareMathOperator*{\argmin}{arg\,min}
\newcommand{\R}{\mathbb{R}}
\newcommand{\SE}{\mathrm{SE}(3)}
\newcommand{\so}{\mathrm{SO}(3)}
\newcommand{\se}{\mathfrak{se}(3)}
\newcommand{\given}{\,\big|\,}
\newcommand{\Bern}{\mathrm{Bern}}
\newcommand{\Norm}{\mathcal{N}}
% (No box-plus / box-minus macros; we use plain Exp / Log throughout.)

\hypersetup{pdftitle={Bernoulli-EKF Multi-Target Tracker on SE(3)}}

\title{Bernoulli-EKF Multi-Target Tracker on
       \texorpdfstring{$\SE$}{SE(3)}\\
       \large A Minimal PMBM Reduction for Manipulation-Scene Object Tracking}
\author{}
\date{}

\begin{document}\maketitle

%-----------------------------------------------------------------------------
\section{Problem formulation}\label{sec:problem}
%-----------------------------------------------------------------------------

\subsection{Setting, frames, and per-object state}\label{subsec:setting}

A calibrated RGB-D camera mounted at a fixed extrinsic
$T_{bc}\!\in\!\SE$ on a mobile robot observes a tabletop workspace
populated by an unknown, time-varying number of rigid objects.
We work with three frames throughout:
\begin{itemize}[leftmargin=2em,topsep=2pt,itemsep=0pt]
\item the world (map) frame $w$,
\item the robot base frame $b$, and
\item the camera frame $c$, with $T_{bc}$ a known kinematic constant
      (read from the URDF; on Fetch-class robots this is a fixed
      head-to-base transform). Camera-in-base lets us interconvert
      $T_{co}=T_{bc}^{-1}T_{bo}$ at zero variance.
\end{itemize}

Two external streams are treated as black boxes:
\begin{itemize}[leftmargin=2em,topsep=0pt]
\item \textbf{Localization} supplies, at every frame $k$, the
      base-in-world Gaussian pose $T_{wb,k}\!\in\!\SE$,
      $\Sigma_{wb,k}\!\in\!\R^{6\times 6}$. Crucially, only its
      \emph{relative} change between consecutive frames,
      \(
         \Delta T_{wb,k} \;\triangleq\; T_{wb,k}\,T_{wb,k-1}^{-1}\!\in\!\SE,
      \)
      enters the recursion as a control input
      (\S\ref{sec:ekf}, eq.~\ref{eq:control}); the absolute
      $(T_{wb,k}, \Sigma_{wb,k})$ enters \emph{only} at the
      world-frame output composition step (\S\ref{sec:emit}). $\sum_{wb,k}$
      \emph{never} appears inside the per-object recursion --- this is
      the key design choice that decouples object-pose tracking from
      localization uncertainty.
\item \textbf{Proprioception} supplies the gripper-in-base pose
      $T_{bg,k}\!\in\!\SE$ from joint encoders (deterministic at the
      Bayesian level). Its relative change
      \(
         \Delta T_{bg,k} \;\triangleq\; T_{bg,k}\,T_{bg,k-1}^{-1}
      \)
      is the control input for held / manipulation-set objects.
\end{itemize}

\paragraph{Per-object state.}
The tracker maintains, online, a variable-size indexed family of
\emph{track parameters}
\begin{equation}
   \mathcal X_k \;=\; \bigl\{X^{(i)}_k\bigr\}_{i\in\mathcal I_k},
   \qquad
   X^{(i)}_k \;=\; \bigl(\,\ell^{(i)},\ T^{(i)}_k,\ P^{(i)}_k,\ r^{(i)}_k\,\bigr),
   \label{eq:obj_state}
\end{equation}
with four attributes per index: a class label $\ell^{(i)}\!\in\!\mathcal L$
(fixed at birth), a \emph{base-frame pose}
$T^{(i)}_k\!\in\!\SE$ (object-in-base, $T_{bo}^{(i)}$ in the long form),
its covariance $P^{(i)}_k\!\in\!\R^{6\times 6}$ in the $\se$ tangent at
$T^{(i)}_k$, and an existence probability $r^{(i)}_k\!\in\![0,1]$ (the
prior-probability that the object referred to by index $i$ actually
exists at time $k$). The index set $\mathcal I_k\subseteq\mathbb{N}$ is
maintained by the tracker (birth enlarges it, pruning shrinks it).

The physical multi-target state at time $k$ is a \emph{random finite
set} $\mathcal T_k\subseteq\SE$ of \emph{base-frame} poses, one per
object actually present. Its (unknown) cardinality is $n_k=|\mathcal T_k|$.
The parameter family $\mathcal X_k$ is the tracker's summary of the
posterior on $\mathcal T_k$ \emph{conditioned on the localization stream}
$\{T_{wb,k},\Sigma_{wb,k}\}_k$; the relationship is made precise in
\S\ref{sec:bern}. World-frame quantities $T_{wo}^{(i)}$ are recovered
from the base-frame state by composition with $T_{wb,k}$ at output
(\S\ref{sec:emit}).

\paragraph{Why base frame.}
The naïve world-frame alternative (store $T_{wo}^{(i)}$, fold
$\Sigma_{wb,k}$ into the measurement noise $R_k$ via
$R_{\mathrm{eff}}=R_{\mathrm{icp}}+\Sigma_{wb,k}$) is correct for one
frame but \emph{double-counts} $\Sigma_{wb,k}$ across consecutive
fusions, because $T_{wb,k}$ and $T_{wb,k+1}$ are nearly the same random
variable. That recursion drives $\Sigma_{wo}$ below $\Sigma_{wb}$,
which is physically impossible for a static object. Storing in base
frame removes $\Sigma_{wb,k}$ from the recursion entirely; it appears
only once, in the world-frame composition at output, where it
correctly lower-bounds $\Sigma_{wo}$.

A per-object shape summary $\mathcal G^{(i)}_k$ (TSDF volume plus
accumulated masked point cloud, rigidly attached to $T^{(i)}_k$) is
maintained as a \emph{by-product} of tracking: fused from depth after
the pose is settled and $r^{(i)}_k$ is high, and fed back into the
recursion only through a predicted-visibility predicate (\S\ref{sec:vis}).
Geometric relations among objects and manipulation-role tags are
deterministic predicates over $\mathcal X_k$, computed after each
recursion step and consumed by downstream modules; they do not
participate in the Bayesian update.

\subsection{Per-frame measurements}

Each frame delivers
\begin{equation}
   Z_k \;=\;\bigl(\,I_k,\ D_k,\ \mathcal D_k\,\bigr),
   \qquad
   \mathcal D_k \;=\; \bigl\{\,(m_{k,l},\ \hat\ell_{k,l},\ \hat s_{k,l},\ \tau_{k,l})\,\bigr\}_{l=1}^{M_k},
   \label{eq:meas_raw}
\end{equation}
where $I_k$ and $D_k$ are RGB and depth images, and each entry of
$\mathcal D_k$ is a binary mask $m_{k,l}\!\in\!\{0,1\}^{H\times W}$ with
class label $\hat\ell_{k,l}\!\in\!\mathcal L$, calibrated detection
score $\hat s_{k,l}\!\in\![0,1]$, and a video-tracklet identifier
$\tau_{k,l}\!\in\!\mathbb N$ issued by the segmenter (SAM2), which
mostly preserves object identity across frames but is occasionally
reset when the segmenter loses track -- it enters the assignment cost
\eqref{eq:cost} as a matching-ID bonus. All four fields are produced
by the open-vocabulary detector + video-segmenter pipeline.
\textbf{Detections live in pixel space; they do not carry 3D pose
information.}

\subsection{The mask-to-pose bottleneck}\label{subsec:bottleneck}

The feature distinguishing this setting from classical multi-target
tracking is that \emph{data association must be resolved before any
pose measurement exists.} The segmenter tells us the image-space extent
of each object hypothesis; converting that to a 3D observation requires
a registration target, and that target is the tracker's own shape model
$\mathcal G^{(i)}_{k-1}$. Each frame therefore executes
\begin{equation}
  \mathcal D_k
  \xrightarrow{\text{back-projection}}
  \bigl\{\mathcal C_{k,l}\bigr\}
  \xrightarrow{\text{association}}
  \pi^{\!*}
  \xrightarrow{\text{ICP vs }\mathcal G^{(i)}_{k-1}}
  \bigl\{(\hat T_{k,l},\,\hat R_{k,l})\bigr\}
  \xrightarrow{\text{filter update}}
  \mathcal X_k.
  \label{eq:pipeline}
\end{equation}
Concretely, for each mask $m_{k,l}$:
\begin{enumerate}[leftmargin=2em,topsep=2pt,itemsep=1pt]
\item \textbf{Back-projection.} The masked pixels of $D_k$ yield a \emph{camera-frame} point cloud $\mathcal C_{k,l}\subset\R^3$ via the intrinsics $K$ and the depth values; localization is \emph{not} needed for this step.
\item \textbf{Association.} A cost-matrix assignment (\S\ref{sec:data-assoc}) decides whether $\mathcal C_{k,l}$ belongs to an existing track $i$, or spawns a new one.
\item \textbf{Pose measurement.} For a matched pair, ICP of $\mathcal C_{k,l}$ against $\mathcal G^{(i)}_{k-1}$ (the latter expressed in the object-local frame and initialised at the prior mean $T^{(i)}_{k|k-1}$ projected to camera space) produces a camera-frame pose estimate $\hat T^{co}_{k,l}\!\in\!\SE$ and a covariance $\hat R^{co}_{k,l}\!\in\!\R^{6\times 6}$ derived from ICP fitness / residual. We then promote both to the base frame using the fixed kinematic $T_{bc}$:
\begin{equation}
   \hat T_{k,l} \;=\; T_{bc}\,\hat T^{co}_{k,l},\qquad
   \hat R_{k,l} \;=\; \operatorname{Ad}(T_{bc})\,\hat R^{co}_{k,l}\,
                       \operatorname{Ad}(T_{bc})^{\!\top}.
   \label{eq:meas_compose}
\end{equation}
$\Sigma_{wb,k}$ is \emph{not} folded into $\hat R_{k,l}$ here -- it
plays no role until output composition (\S\ref{sec:emit}). For an
unmatched mask we take $\hat T^{co}_{k,l}$ from the centroid of
$\mathcal C_{k,l}$ with identity rotation and inflate $\hat R^{co}_{k,l}$
as a birth prior.
\end{enumerate}
Only after \eqref{eq:pipeline} runs does the filter consume
$(\hat T_{k,l},\hat R_{k,l})$ as the measurement of an EKF update
(\S\ref{sec:ekf}). The tracker is therefore a feedback loop in which
yesterday's state is what makes today's measurement well-defined.

\subsection{Estimation target and pathologies}

Conditioned on the localization and proprioception streams, the
Bayesian target is the multi-target set posterior over
\emph{base-frame} poses
\begin{equation}
   f_k(\mathcal T) \;\triangleq\;
       p\!\Bigl(\mathcal T_k=\mathcal T \;\Big|\;
                Z_{1:k},\,T_{wb,1:k},\,\Sigma_{wb,1:k},\,T_{bg,1:k}\Bigr).
   \label{eq:target}
\end{equation}
The tracker represents $f_k$ by the family $\mathcal X_k$ of
\eqref{eq:obj_state} under the Multi-Bernoulli assumption of
\S\ref{sec:bern}. Downstream we report the maximum-likelihood cardinality
$\hat n_k$, per-track labels $\ell^{(i)}$, pose Gaussians
$(T^{(i)}_k,P^{(i)}_k)$, and existence probabilities $r^{(i)}_k$.

The recursion must absorb four failure modes simultaneously:
\begin{enumerate}[leftmargin=2em,label=(\roman*),topsep=2pt,itemsep=0pt]
\item\label{p:miss}     \emph{missed detections} -- an object is present but yields no mask, either because the detector failed (probability $1-p_d$) or because it is occluded / out-of-frustum (visibility $p_v$);
\item\label{p:clutter}  \emph{clutter} -- spurious masks not corresponding to any real object, at intensity $\lambda_c$;
\item\label{p:assoc}    \emph{association ambiguity} -- which mask belongs to which track; a wrong answer corrupts \emph{both} the pose belief (wrong measurement) \emph{and} the existence belief (spurious confirmation);
\item\label{p:mismatch} \emph{model mismatch} -- pose changes exceed what $P^{(i)}_{k|k-1}$ predicts, e.g.\ an object displaced by an unmodelled external agent, or ICP converging to a wrong local minimum.
\end{enumerate}
Data association \ref{p:assoc} is one failure mode of four. The
construction below treats it as a single pass through the same
Mahalanobis-distance machinery that gates \ref{p:clutter} /
\ref{p:mismatch} and weights \ref{p:miss}: a single scalar statistic
$d^2$ drives every decision.

We formalise the recursion as a \emph{Bernoulli-EKF tracker} on $\SE$:
per-object Gaussian pose beliefs with explicit existence probabilities,
updated via a single-hypothesis Hungarian assignment on the $d^2$ cost
matrix, with birth from unmatched measurements and principled pruning
from the existence update. This is the small-cardinality, low-clutter,
single-hypothesis limit of the Poisson Multi-Bernoulli Mixture
recursion\,\cite{Williams2015,GarciaFernandez2018} -- equivalently, the
JIPDA reduction on $\SE$.

%-----------------------------------------------------------------------------
\section{$\SE$ preliminaries}\label{sec:se3}
%-----------------------------------------------------------------------------

We treat poses as elements of the matrix Lie group $\SE$ and represent
small pose perturbations as 6-vectors $\xi=(\rho,\omega)\!\in\!\R^6$,
collecting a translation $\rho\!\in\!\R^3$ and a rotation vector
$\omega\!\in\!\R^3$. The mapping $\xi\!\mapsto\!\xi^{\wedge}\!\in\!\se$
to the Lie algebra $\se$ is the standard one\,\cite{MurrayLiSastry1994};
its inverse is denoted $\xi=(\xi^{\wedge})^{\vee}$. We invoke
$\xi^{\wedge}$ only to state two basic operations in closed form, after
which all computation happens directly in $\R^6$:
\begin{align}
   \Exp \,:\, \R^6 \to \SE, &\qquad
      \Exp(\xi) \;\triangleq\; \exp(\xi^{\wedge}),
      \label{eq:expmap}\\
   \Log \,:\, \SE \to \R^6, &\qquad
      \Log(T)  \;\triangleq\; \log(T)^{\vee}.
      \label{eq:logmap}
\end{align}
$\Exp$ has a closed form (Rodrigues' formula in the rotation block, with
a translation Jacobian)\,\cite{MurrayLiSastry1994}; $\Log$ is its
inverse. They are the only structural operations on $\SE$ we use; group
composition is plain matrix multiplication and inversion is the matrix
inverse.

A \emph{Gaussian on $\SE$} centred at $T_0\!\in\!\SE$ with covariance
$P\!\in\!\R^{6\times6}$ is the distribution of
\begin{equation}
   T \;=\; T_0\,\Exp(\xi),\qquad \xi\sim\Norm(\bm 0,P),
   \label{eq:gauss_se3}
\end{equation}
written $T\sim\Norm(T_0,P)$. The mean lives in $\SE$; the uncertainty
lives in $\R^6$, attached to the tangent space at $T_0$.

%-----------------------------------------------------------------------------
\section{Extended Kalman filter on $\SE$ (single target)}\label{sec:ekf}
%-----------------------------------------------------------------------------

\subsection*{Overview and notation}

Treat each tracked object in isolation. The pose belief at time $k$ is a
Gaussian on $\SE$ in the body-tangent parametrisation of \S\ref{sec:se3}:
the mean is $T_k\!\in\!\SE$ (a $4\!\times\!4$ rigid transform of the
object frame in the world frame), and the spread around it is the
$6\!\times\!6$ covariance
\(
   P_k \;=\; \mathrm{cov}\bigl(\xi_k\bigr),\;\; \xi_k\sim\Norm(\bm 0, P_k),
\)
with $\xi=(\rho,\omega)$ the local twist that perturbs $T_k$ via
$T_k\,\Exp(\xi)$. When we write $T_{k|k-1}, P_{k|k-1}$ we mean the
\emph{predicted} mean / covariance after the process step but before
incorporating the frame-$k$ measurement; $T_{k|k}, P_{k|k}$ are their
\emph{posterior} counterparts after the update.

The recursion is the standard Extended Kalman filter on a Lie group,
specialised to (a) constant-pose dynamics and (b) right-multiplicative,
tangent-Gaussian process / measurement noise. Each frame:

\begin{enumerate}[leftmargin=2em,topsep=2pt,itemsep=1pt]
\item \textbf{Predict.} Apply the known control input $u_k\!\in\!\SE$
      (\S\ref{subsec:setting}: $\Delta T_{wb,k}^{-1}$ for world-static
      objects, $\Delta T_{bg,k}$ for held / manipulation-set objects,
      $u_k=I$ as the trivial fallback) to the mean by left
      multiplication, $\operatorname{Ad}$-conjugate the covariance, add
      $Q_k$, then cap with the saturation projector $\Phi$. $Q_k$
      encodes the residual per-frame dynamics not explained by $u_k$:
      tiny for a stationary world-static object (most of $Q_k$ work is
      done by $u_k=\Delta T_{wb,k}^{-1}$), small slack
      $Q_{\text{manip}}$ for held objects (grip slip / deformation),
      large during grasp/release transitions when $u_k$ is unreliable.
      Supplied by the manipulation context
      (\texttt{ekf\_se3.process\_noise\_for\_phase}).
\item \textbf{Measure.} The frame's mask is back-projected to a
      \emph{camera-frame} 3-D cloud and registered (ICP,
      \S\ref{subsec:bottleneck}) against the track's shape summary,
      yielding $\hat T^{co}_{k,l}\!\in\!\SE$ and
      $\hat R^{co}_{k,l}\!\in\!\R^{6\times 6}$. The fixed kinematic
      $T_{bc}$ promotes both to the base frame at zero variance:
      $z_k=T_{bc}\hat T^{co}_{k,l}$,
      $R_k=\operatorname{Ad}(T_{bc})\hat R^{co}_{k,l}\operatorname{Ad}(T_{bc})^{\!\top}$
      (eq.~\ref{eq:meas_compose}). $\Sigma_{wb,k}$ does \emph{not}
      enter $R_k$.
\item \textbf{Update.} Linearise both the prior Gaussian and the
      observation in the tangent at $T_{k|k-1}$. The linearised
      innovation $\nu_k=\Log\!\bigl(T_{k|k-1}^{-1}z_k\bigr)\!\in\!\R^6$
      lives in the same $6$-vector space as $\xi$; with $H_k=I_6$ the
      Kalman gain is $K_k = P_{k|k-1}S_k^{-1}$, $S_k=P_{k|k-1}+R_k$.
      The mean is corrected on the manifold via right-multiplication by
      $\Exp(K_k\nu_k)$; the covariance is updated in Joseph form and
      then $\Phi$-saturated.
\end{enumerate}

\paragraph{Glossary of variables (this section).}
\begin{description}[leftmargin=2em,topsep=2pt,itemsep=0pt,labelindent=0em]
\item[$T_k$, $P_k$.] Posterior \emph{base-frame} object pose mean and
      tangent covariance at frame $k$ ($T_{bo,k}^{(i)}$ in the long
      form). $T_k\!\in\!\SE$ as a $4\!\times\!4$ matrix;
      $P_k\!\in\!\R^{6\times 6}$ symmetric PD. Maintained recursively
      (initialised at birth, \S\ref{sec:birth}; never measured directly).
\item[$T_{k|k-1}$, $P_{k|k-1}$.] Predicted mean / covariance produced by step~1.
\item[$u_k$.] Known control input in $\SE$ supplied by external streams:
      \[
        u_k =
        \begin{cases}
          \Delta T_{bg,k} = T_{bg,k}T_{bg,k-1}^{-1}, & i\in\text{manipulation set} \\
          \Delta T_{wb,k}^{-1} = T_{wb,k-1}T_{wb,k}^{-1}, & \text{otherwise (world-static).}
        \end{cases}
      \]
      The first form is the gripper change in base (proprioception);
      the second is the inverse base motion (a world-static object's
      base-frame pose moves opposite to the base). $u_k = I$ on the
      first frame (no $\Delta$).
\item[$\xi_k$, $\nu_k$.] $\R^6$ vectors in the body tangent at
      $T_{k-1}$ and $T_{k|k-1}$ respectively. $\xi_k$ is the
      \emph{unmodelled} per-frame perturbation (anything not captured by
      $u_k$); $\nu_k$ is the observed innovation.
\item[$Q_k$.] $6\!\times\!6$ process-noise covariance for the residual
      after $u_k$ is applied, supplied externally per manipulation
      phase (\texttt{process\_noise\_for\_phase}; not estimated).
\item[$z_k$, $R_k$.] \emph{Base-frame} measurement and its covariance.
      $z_k = T_{bc}\hat T^{co}_{k,l}\!\in\!\SE$ from per-mask ICP
      against $\mathcal G^{(i)}_{k-1}$ promoted to base via $T_{bc}$;
      $R_k = \operatorname{Ad}(T_{bc})\hat R^{co}_{k,l}\operatorname{Ad}(T_{bc})^{\!\top}$
      blends ICP residuals only -- $\Sigma_{wb,k}$ is excluded.
\item[$T_{bc}$.] Fixed camera-in-base extrinsic (kinematic constant).
\item[$\Delta T_{wb,k}$, $\Delta T_{bg,k}$.] Frame-to-frame deltas of
      base-in-world (from localization) and gripper-in-base (from
      proprioception); see \S\ref{subsec:setting}.
\item[$S_k$, $K_k$, $H_k$.] Innovation covariance, Kalman gain,
      observation Jacobian. $H_k=I_6$ to first order in $\nu_k$.
\item[$\Phi$, $P_{\max}$.] Covariance-saturation projector
      (eq.~\ref{eq:phi}) and its design constant $P_{\max}\!\succ\!\bm 0$;
      caps $\Tr P$ at $\Tr P_{\max}$ so the recursion is self-bounded
      under arbitrarily long missing-observation streaks.
\item[$d_k^2$, $\mathcal L(z_k\!\mid\!\cdot)$.] Mahalanobis squared
      distance $\nu_k^{\!\top}S_k^{-1}\nu_k$ (eq.~\ref{eq:mahal}) and
      the Gaussian innovation likelihood (eq.~\ref{eq:ekf_lik}); both
      consumed downstream by the Bernoulli existence update
      (\S\ref{sec:update}) and the Hungarian cost matrix
      (\S\ref{sec:data-assoc}).
\end{description}

The remainder of the section makes the three steps quantitative.

\subsection*{Process and measurement models}

Tracking happens in the \emph{robot base frame} throughout. For one
rigid object, the process model is a known left-multiplicative control
$u_k\!\in\!\SE$ followed by $\R^6$ tangent-space process noise applied
on the right:
\begin{equation}
   T_k \;=\; u_k\,T_{k-1}\,\Exp(\xi_k),\qquad
   \xi_k\sim\Norm(\bm 0,Q_k).\label{eq:proc}
\end{equation}
$u_k$ encodes how a deterministic external action moves the object's
\emph{base-frame} pose between frames:
\begin{equation}
   u_k \;=\;
   \begin{cases}
      T_{bg,k}\,T_{bg,k-1}^{-1},
        & i\in\text{manipulation set (held / on / in)}, \\[2pt]
      T_{wb,k}^{-1}\,T_{wb,k-1},
        & \text{otherwise (world-static object)}.
   \end{cases}
   \label{eq:control}
\end{equation}
Both expressions are derived below from the constraint that the
\emph{world-frame} pose of the object is conserved (static) or
identified with the gripper (held). In both cases $u_k$ is the
\textbf{body-frame} inverse of the base/gripper motion, not the
world-frame inverse: $T_{wb,k}^{-1}T_{wb,k-1}\!\ne\!T_{wb,k-1}T_{wb,k}^{-1}$
in general. Mistaking the latter for the former produces $\sim$1\,m
drift on a typical mobile-manipulator trajectory because the conjugate
transformation differs by an adjoint rotation -- enforced experimentally
on the apple\_in\_the\_tray sequence (\S\ref{sec:approx} A.7 caveat).

Both $T_{wb,k}$ and $T_{bg,k}$ are external streams
(\S\ref{subsec:setting}); we treat them as point estimates -- the
implementation feeds in $T_{wb,k}$ as the SLAM mean and $T_{bg,k}$ from
joint encoders. Localization uncertainty $\Sigma_{wb,k}$ does
\emph{not} enter $u_k$ or $Q_k$; it appears once at the world-frame
output composition (\S\ref{sec:emit}).

\paragraph{Justification of the two regimes.}
\emph{World-static object.} If $T_{wo}^{(i)}$ is constant then
\(
   T_{wb,k}\,T_{bo,k}^{(i)} = T_{wb,k-1}\,T_{bo,k-1}^{(i)}
\)
and so
\(
   T_{bo,k}^{(i)}
     = T_{wb,k}^{-1}\,T_{wb,k-1}\,T_{bo,k-1}^{(i)}
     = u_k\,T_{bo,k-1}^{(i)}
\)
with $u_k=T_{wb,k}^{-1}T_{wb,k-1}$.
\emph{Held object.} If the object is rigidly attached to the gripper
($T_{go}^{(i)}\!=$const), then
\(
   T_{bo,k}^{(i)}=T_{bg,k}T_{go}^{(i)}
     =T_{bg,k}T_{bg,k-1}^{-1}T_{bo,k-1}^{(i)}
     =u_k\,T_{bo,k-1}^{(i)}
\)
with $u_k=T_{bg,k}T_{bg,k-1}^{-1}$ -- the gripper's relative motion is
naturally a left-form transform here because $T_{bo}$ and $T_{bg}$
share the base frame on the left of the composition. Setting $u_k=I$
recovers the standard constant-pose model.

$Q_k$ encodes the residual per-frame dynamics not explained by $u_k$:
small slack for held objects (grip slip), tiny drift for static objects
(SLAM noise leaks already captured by $u_k$), large during
grasp/release transitions when the regime is switching and $u_k$ is
unreliable. Supplied externally by the manipulation context.

The measurement is an ICP-derived noisy pose lifted to the base frame
(cf.\ \S\ref{subsec:bottleneck}), modelled the same way:
\begin{equation}
   z_k \;=\; T_k\,\Exp(\eta_k),\qquad \eta_k\sim\Norm(\bm 0,R_k),
   \label{eq:meas}
\end{equation}
where $z_k = T_{bc}\hat T^{co}_{k,l}$ and
$R_k = \operatorname{Ad}(T_{bc})\hat R^{co}_{k,l}\operatorname{Ad}(T_{bc})^{\!\top}$
(eq.~\ref{eq:meas_compose}). $R_k$ absorbs only ICP uncertainty;
$\Sigma_{wb,k}$ is excluded by construction.

\paragraph{Predict.} Eq.\,\eqref{eq:proc} composes a known SE(3)
left-multiplication $u_k$ with right-multiplicative tangent noise.
Writing the prior as $T_{k-1|k-1}\Exp(\xi_{k-1|k-1})$ with
$\xi_{k-1|k-1}\!\sim\!\Norm(\bm 0,P_{k-1|k-1})$, the predicted pose is
\begin{align*}
T_k &= u_k\,T_{k-1|k-1}\,\Exp(\xi_{k-1|k-1})\,\Exp(\xi_k) \\
    &= \bigl[u_k\,T_{k-1|k-1}\bigr]\,
       \Exp\!\bigl(\xi_{k-1|k-1}+\xi_k+\tfrac12[\xi_{k-1|k-1},\xi_k]+\cdots\bigr)
\end{align*}
by Baker--Campbell--Hausdorff. Dropping second-order and higher
brackets (the standard EKF-on-Lie-group linearisation) yields
$\xi_{k|k-1}\!\approx\!\xi_{k-1|k-1}+\xi_k\!\sim\!\Norm(\bm 0,
P_{k-1|k-1}+Q_k)$ in the tangent at the \emph{shifted} mean
$u_kT_{k-1|k-1}$. Because $u_k$ acts on the left of the original mean,
the original tangent vectors transport to the new mean by the adjoint
$\operatorname{Ad}(u_k)$\,\cite{MurrayLiSastry1994}, so the prior
covariance reaches the new tangent as
$\operatorname{Ad}(u_k)\,P_{k-1|k-1}\,\operatorname{Ad}(u_k)^\top$.
Applying the \emph{covariance-saturation projector}
\begin{equation}
   \Phi:\R^{6\times6}\to\R^{6\times6},\qquad
   \Phi(P) \;\triangleq\; \min\!\bigl(1,\,\Tr P_{\max}/\Tr P\bigr)\,P,
   \label{eq:phi}
\end{equation}
for a chosen design constant $P_{\max}\!\succ\!\bm 0$ (a static
uninformative prior, $\Tr P_{\max}\!>\!0$), the predict step is
\begin{align}
   T_{k|k-1} &= u_k\,T_{k-1|k-1},\label{eq:ekf_T_pred}\\
   P_{k|k-1} &= \Phi\!\bigl(\operatorname{Ad}(u_k)\,P_{k-1|k-1}\,
                              \operatorname{Ad}(u_k)^\top + Q_k\bigr).
                \label{eq:ekf_P_pred}
\end{align}
When $u_k\!=\!I$ this collapses to the constant-pose form $T_{k|k-1}=T_{k-1|k-1}$,
$P_{k|k-1}=\Phi(P_{k-1|k-1}+Q_k)$. $\Phi$ preserves symmetry and
positive-definiteness (scaling by a positive scalar) and caps the trace
at $\Tr P_{\max}$: the recursion is self-bounded under arbitrarily long
missing-observation streaks. Without it $P$ grows linearly in the
number of missed frames.

\paragraph{Update.} The innovation, residual covariance, and Kalman gain
are
\begin{equation}
   \nu_k \;=\; \Log\!\bigl(T_{k|k-1}^{-1}\,z_k\bigr) \;\in\;\R^6,\qquad
   S_k   \;=\; P_{k|k-1} + R_k,\qquad
   K_k   \;=\; P_{k|k-1}\,S_k^{-1}.\label{eq:ekf_innov}
\end{equation}
The measurement Jacobian is $H_k\!=\!I_6$ to first order in $\nu_k$:
both $z$ and $T$ live in $\SE$, and the $\Log$ is evaluated at the prior
mean, so the linearisation is exact when $\nu_k\!=\!\bm 0$ and degrades
smoothly with $\|\nu_k\|$. The state correction is applied in the
tangent and folded back via right-multiplication by $\Exp$; covariance
is updated in Joseph form for numerical stability:
\begin{align}
   T_{k|k} &= T_{k|k-1}\,\Exp\!\bigl(K_k\,\nu_k\bigr),\label{eq:ekf_T_upd}\\
   P_{k|k} &= \Phi\!\Bigl(\,
                (I_6-K_k)P_{k|k-1}(I_6-K_k)^\top + K_kR_kK_k^\top
                \Bigr).\label{eq:ekf_P_upd}
\end{align}
The same projector $\Phi$ is applied after the update so a
high-innovation observation cannot inflate $P_{k|k}$ above $P_{\max}$
either.

\paragraph{Innovation likelihood and Mahalanobis statistic.}
By \eqref{eq:meas} and the linearisation of $\Log$ at $T_{k|k-1}^{-1}T_k$,
the innovation satisfies $\nu_k\sim\Norm(\bm 0, S_k)$ under the null
hypothesis $H_0^{(k)}$: ``$z_k$ was generated by the track whose prior
is $\Norm(T_{k|k-1},P_{k|k-1})$.'' The corresponding Gaussian density
on $\nu_k\in\R^6$, which we use as the measurement likelihood for the
Bernoulli update below, is
\begin{equation}
   \mathcal L(z_k\given T_{k|k-1},P_{k|k-1})
   \;\triangleq\;
   \Norm\!\bigl(\nu_k;\,\bm 0,\,S_k\bigr)
   \;=\;
   (2\pi)^{-3}|S_k|^{-1/2}\exp\!\bigl(-\tfrac12 d_k^2\bigr),
   \label{eq:ekf_lik}
\end{equation}
\begin{equation}
   d_k^2 \;\triangleq\; \nu_k^\top S_k^{-1}\nu_k \;\sim\; \chi^2_6\ \text{under}\ H_0^{(k)}.
   \label{eq:mahal}
\end{equation}
\eqref{eq:ekf_lik} is exact at $\nu_k\!=\!\bm 0$ (where the Jacobian
of $z\!\mapsto\!\Log(T_{k|k-1}^{-1}z)$ is $I_6$) and first-order accurate
elsewhere; the tangent density agrees with the $\SE$ density up to this
Jacobian.

\paragraph{Robust update.} Define an inner gate
$G_{\text{in}}=\chi^2_6(0.95)\approx 12.59$ and an outer gate
$G_{\text{out}}=\chi^2_6(0.9997)\approx 25$, and the re-weighting factor
\begin{equation}
   w_k \;=\;
   \begin{cases}
      1,                                 & d_k^2 \le G_{\text{in}},\\
      \sqrt{G_{\text{in}}/d_k^2},        & G_{\text{in}} < d_k^2 \le G_{\text{out}},\\
      0,                                 & d_k^2 > G_{\text{out}}.
   \end{cases}\label{eq:huber}
\end{equation}
The update \eqref{eq:ekf_innov}--\eqref{eq:ekf_P_upd} is then modified as follows:
\begin{itemize}[leftmargin=2em,topsep=0pt,itemsep=0pt]
\item If $w_k\!=\!0$ (outer-gate reject): skip the update entirely,
      setting $T_{k|k}\!=\!T_{k|k-1}$, $P_{k|k}\!=\!P_{k|k-1}$; $z_k$
      is routed to the unassigned-track branch of the Bernoulli update
      (\S\ref{sec:update}).
\item Otherwise ($w_k\!>\!0$): substitute $R_k\!\leftarrow\!R_k/w_k$ in
      \eqref{eq:ekf_innov}--\eqref{eq:ekf_P_upd}. Note that
      $K_k,S_k$ are recomputed with the inflated $R_k/w_k$; $\nu_k$ is
      unchanged.
\end{itemize}
\eqref{eq:huber} is Huber's redescending M-estimator on the Mahalanobis
residual, applied to the re-weighted EKF.

%-----------------------------------------------------------------------------
\section{Per-object Bernoulli state}\label{sec:bern}
%-----------------------------------------------------------------------------

To each index $i\in\mathcal I_k$ we associate a \emph{Bernoulli
component} with parameters $B^{(i)} = (r^{(i)},\ell^{(i)},T^{(i)},P^{(i)})$,
representing the single-target set posterior over a per-index random
finite set $\mathcal T^{(i)}\subseteq\SE$ with set density
\begin{equation}
   f^{(i)}(\mathcal T) \;=\;
   \begin{cases}
       1-r^{(i)},                                   & \mathcal T = \varnothing,\\
       r^{(i)}\,\Norm\!\bigl(T;\,T^{(i)},P^{(i)}\bigr), & \mathcal T = \{T\},\ T\!\in\!\SE,\\
       0,                                           & |\mathcal T|\ge 2.
   \end{cases}\label{eq:bern_density}
\end{equation}
Here $\Norm(\cdot;T^{(i)},P^{(i)})$ is the Gaussian density on $\SE$ of
\eqref{eq:gauss_se3}. In words: with probability $1-r^{(i)}$ the $i$-th
track does not exist (empty set); with probability $r^{(i)}$ it exists
at exactly one pose $T\!\in\!\SE$ distributed according to
$\Norm(\cdot;T^{(i)},P^{(i)})$. Classical Bernoulli-RFS normalisation
$\int f^{(i)}(\mathcal T)\,\delta\mathcal T = 1$ holds by construction
(FISST set integral; see \cite{Williams2015}). The pair $(T^{(i)},P^{(i)})$
is the EKF state of \S\ref{sec:ekf}; $r^{(i)}$ is the new variable PMBM
contributes; $\ell^{(i)}$ is fixed at birth and does not enter the RFS
marginalisation.

\paragraph{Multi-Bernoulli approximation.} We assume the per-index
Bernoulli sets $\mathcal T^{(1)},\dots,\mathcal T^{(n)}$ are
statistically independent. The physical pose RFS of \S\ref{sec:problem}
is then
\begin{equation}
   \mathcal T_k \;=\; \bigcup_{i\in\mathcal I_k}\mathcal T^{(i)}_k,
   \label{eq:pose_rfs_union}
\end{equation}
and its multi-target set density factorises as the convolution (FISST
union) of the per-index Bernoulli densities. Equivalently, for any
measurable subset $\mathcal T\!\subseteq\!\SE$,
\begin{equation}
   f_k(\mathcal T)\ =\!\!\!\!\sum_{\substack{\{\mathcal T^{(i)}\}_{i\in\mathcal I_k}:\\ \bigcup_i\mathcal T^{(i)} = \mathcal T}}
        \;\prod_{i\in\mathcal I_k} f^{(i)}\!\bigl(\mathcal T^{(i)}\bigr),
   \label{eq:mb_factorisation}
\end{equation}
which is the Multi-Bernoulli approximation to \eqref{eq:target}.
Independence is the only approximation introduced relative to a full
PMBM mixture; the class labels $\{\ell^{(i)}\}$ ride along as
track-identity metadata.

%-----------------------------------------------------------------------------
\section{Bernoulli prediction}\label{sec:bern-pred}
%-----------------------------------------------------------------------------

For each $i\in\mathcal I_{k-1}$, given a constant survival probability
$p_s\!\in\!(0,1]$, prediction is the standard PMBM step (eq.\,(43) in
\cite{GarciaFernandez2018}, simplified to constant $p_s$):
\begin{align}
   (T^{(i)}_{k|k-1},P^{(i)}_{k|k-1})
       &\quad\text{via EKF predict
                  \eqref{eq:ekf_T_pred}--\eqref{eq:ekf_P_pred} with
                  $u^{(i)}_k$, $Q^{(i)}_k$,}\nonumber\\[2pt]
   r^{(i)}_{k|k-1} &= p_s\,r^{(i)}_{k-1|k-1}.\label{eq:bern_pred_r}
\end{align}
$p_s$ is the (observer-independent) physical survival probability of
the object between consecutive frames. For static manipulation scenes,
where objects neither move nor vanish on their own, the natural value
is $p_s\!=\!1$, which collapses \eqref{eq:bern_pred_r} to identity on
$r$. Setting $p_s\!<\!1$ would model spontaneous physical
disappearance (rare in our domain). Crucially, $p_s$ is \emph{not}
visibility-dependent: occlusion-induced ``invisibility'' does not
affect physical survival; it only affects the observer's ability to
detect, which enters the recursion through $\tilde p_d$ in the missed
update (\S\ref{sec:update}). $u^{(i)}_k$ is the per-track control input
(eq.~\ref{eq:control}) that depends on whether $i$ is in the
manipulation set; $Q^{(i)}_k$ is the residual process noise supplied by
the manipulation context.

Object birth is folded into data association (\S\ref{sec:birth}); we do
not maintain a continuous Poisson birth intensity.

%-----------------------------------------------------------------------------
\section{Data association}\label{sec:data-assoc}
%-----------------------------------------------------------------------------

For every track--detection pair $(i,l)\!\in\!\mathcal I_{k-1}\!\times\!\{1,\dots,M_k\}$ compute
\begin{equation}
   \nu_{il} = \Log\!\bigl(T^{(i)\,-1}_{k|k-1}\hat T_{k,l}\bigr),\quad
   S_{il}   = P^{(i)}_{k|k-1} + \hat R_{k,l},\quad
   d^2_{il} = \nu_{il}^\top S_{il}^{-1}\nu_{il},
   \label{eq:innov_pair}
\end{equation}
i.e.\ \eqref{eq:ekf_innov}--\eqref{eq:mahal} with
$(z_k,R_k)\!\leftarrow\!(\hat T_{k,l},\hat R_{k,l})$ and
$(T_{k|k-1},P_{k|k-1})\!\leftarrow\!(T^{(i)}_{k|k-1},P^{(i)}_{k|k-1})$.
The assignment cost matrix
$C\!\in\!(\R\cup\{+\infty\})^{|\mathcal I_{k-1}|\times M_k}$ combines
the Mahalanobis distance with a class-label feasibility gate, a
$\chi^2_6$ outer gate, and -- when the segmenter exposes a tracklet
identifier $\tau_{k,l}$ -- an additive bonus on matching IDs:
\begin{equation}
   C_{il} \;=\;
   \begin{cases}
      d^2_{il} \;-\; \alpha\,\mathbb 1\!\bigl[\tau_{k,l}=\tau^{(i)}_{k-1}\bigr],
        & \hat\ell_{k,l}=\ell^{(i)}\ \text{and}\ d^2_{il}\le G_{\text{out}},\\
      +\infty, & \text{otherwise.}
   \end{cases}
   \label{eq:cost}
\end{equation}
Here $\tau^{(i)}_{k-1}\!\in\!\mathbb N\cup\{\varnothing\}$ is the
tracklet ID with which track $i$ was last matched ($\varnothing$ if
never matched, in which case the indicator is $0$), and
$\alpha\!\ge\!0$ is a tunable bonus in $\chi^2_6$ units ($\alpha\!=\!0$
drops the ID term; the paper default $\alpha\!=\!4.4$ breaks ties
between two same-class tracks whose $d^2$ differ by a few units
without overriding a clear geometric mismatch).

The MAP assignment is
\begin{equation}
   \pi^{\!*} \;=\;
   \argmin_{\pi}\;
   \sum_{i\in\mathcal I_{k-1}\,:\,\pi(i)>0} C_{i,\pi(i)},\qquad
   \text{s.t.}\ \pi(i)=\pi(j)>0\Rightarrow i=j,
   \label{eq:hung}
\end{equation}
with $\pi:\mathcal I_{k-1}\to\{0,1,\dots,M_k\}$ and $\pi(i)\!=\!0$
encoding ``track $i$ unassigned''. \eqref{eq:hung} is solved with
scipy's \texttt{linear\_sum\_assignment} (Jonker--Volgenant); any
cost-$+\infty$ entry forces $\pi^{\!*}(i)\!=\!0$. After solving, set
$\tau^{(i)}_k\!\leftarrow\!\tau_{k,l}$ on matched pairs and leave
$\tau^{(i)}_k\!=\!\tau^{(i)}_{k-1}$ on misses (so a transient tracklet
drop does not erase the memory); birth initialises
$\tau^{(i_{\mathrm{new}})}_k\!=\!\tau_{k,l}$. Measurements with no
positive-assignment track are available for birth (\S\ref{sec:birth}).
\eqref{eq:hung} is the MAP collapse of the PMBM mixture over global
associations (JIPDA reduction).

Note on circularity: $d^2_{il}$ depends on
$(\hat T_{k,l},\hat R_{k,l})$, which is itself produced by ICP against
$\mathcal G^{(i)}_{k-1}$. The orchestrator breaks the loop by running
ICP once per track on a coarsely-gated candidate set (masks whose
world-frame centroid is within a fixed window of $T^{(i)}_{k|k-1}$);
\eqref{eq:cost} then refines the selection using the resulting pose
measurement. Masks failing the coarse gate for all tracks go to birth
with a centroid-initialised $\hat T_{k,l}$.

%-----------------------------------------------------------------------------
\section{Bernoulli measurement update}\label{sec:update}
%-----------------------------------------------------------------------------

The updates below are direct Bayes applications conditioned on the MAP
association $\pi^{\!*}$. Let $E^{(i)}$ denote the event ``track $i$
exists,'' with prior $\Pr[E^{(i)}] = r^{(i)}_{k|k-1}$.

\paragraph{Track $i$ associated with $z_{k,l}$ ($\pi^{\!*}(i)=l$).}
Given the MAP assignment, $z_{k,l}$ is drawn either from the track's
Gaussian-on-$\SE$ (if $E^{(i)}$, detected with prob. $p_d$) or from
clutter (if $\neg E^{(i)}$, density $\lambda_c$). By Bayes,
\begin{align}
   \Pr[E^{(i)} \given z_{k,l}\!\sim\!\text{track}\,i]
     &= \frac{\Pr[z_{k,l}\given E^{(i)}]\,\Pr[E^{(i)}]}
             {\Pr[z_{k,l}\given E^{(i)}]\Pr[E^{(i)}]
              + \Pr[z_{k,l}\given \neg E^{(i)}]\Pr[\neg E^{(i)}]}\nonumber\\
     &= \frac{p_d\,\mathcal L_{il}\,r^{(i)}_{k|k-1}}
             {p_d\,\mathcal L_{il}\,r^{(i)}_{k|k-1}
              + \lambda_c\,\bigl(1-r^{(i)}_{k|k-1}\bigr)},\label{eq:r_assoc_deriv}
\end{align}
where $\mathcal L_{il}$ is \eqref{eq:ekf_lik} evaluated at $(i,l)$.
Hence
\begin{equation}
\boxed{\quad
   r^{(i)}_{k|k} \;=\;
   \frac{p_d\,\mathcal L_{il}\,r^{(i)}_{k|k-1}}
        {p_d\,\mathcal L_{il}\,r^{(i)}_{k|k-1}
         + \lambda_c\,\bigl(1-r^{(i)}_{k|k-1}\bigr)}.
   \quad}\label{eq:r_assoc}
\end{equation}
The conditional state posterior $p(T^{(i)}\given E^{(i)},z_{k,l})$ is
the EKF update of \eqref{eq:ekf_T_upd}--\eqref{eq:ekf_P_upd} with
measurement $(z_k,R_k)\!=\!(\hat T_{k,l},\hat R_{k,l}/w_{il})$ and Huber
weight $w_{il}$ from \eqref{eq:huber}; if $w_{il}\!=\!0$ the pair is
rerouted to the unassigned branch below.

\paragraph{Track $i$ unassigned ($\pi^{\!*}(i)=0$).}
Define the \emph{state-dependent} detection probability
$\tilde p_d^{(i)}\!=\!p_d\cdot p_v^{(i)}$, where $p_v^{(i)}$ is the
geometric visibility predicate of \S\ref{sec:vis} evaluated at
$T^{(i)}_{k|k-1}$. Under the standard PMBM measurement model this is a
state-dependent $p_d(\cdot)$ in PMBM's notation
(eq.\,(2) in \cite{GarciaFernandez2018}); the visibility just makes
the dependence explicit. Absence of a track-$i$ measurement is
consistent with existence only through the missed-detection mechanism
(probability $1-\tilde p_d^{(i)}$); non-existence yields the event
trivially. By Bayes on the event $E^{(i)}\!=\!\text{``track $i$ exists''}$,
\begin{equation}
\boxed{\quad
   r^{(i)}_{k|k} \;=\;
   \frac{(1-\tilde p_d^{(i)})\,r^{(i)}_{k|k-1}}
        {(1-\tilde p_d^{(i)})\,r^{(i)}_{k|k-1} + \bigl(1-r^{(i)}_{k|k-1}\bigr)},
   \quad}\label{eq:r_miss}
\end{equation}
which is the standard Bernoulli-misdetection update with
state-dependent $p_d$ (cf.\ García-Fern\'{a}ndez et al.\ eq.\,(31)).
The state $(T^{(i)},P^{(i)})$ inherits from the prediction since no
new information arrives:
$T^{(i)}_{k|k}\!=\!T^{(i)}_{k|k-1}$ and
$P^{(i)}_{k|k}\!=\!P^{(i)}_{k|k-1}$.

When $p_v^{(i)}\!\to\!0$ (geometric occlusion or out-of-FOV) we have
$\tilde p_d^{(i)}\!\to\!0$, yielding $r^{(i)}_{k|k}\!\to\!r^{(i)}_{k|k-1}$
in \eqref{eq:r_miss}: \emph{no penalty for not seeing the unseen}.
Long-term existence loss for genuinely missing-but-expected tracks
comes from accumulating \eqref{eq:r_miss} over frames in which
$p_v^{(i)}$ is high.

%-----------------------------------------------------------------------------
\section{Birth: unmatched measurements}\label{sec:birth}
%-----------------------------------------------------------------------------

Each measurement index $l$ with $\pi^{\!*}(i)\!\ne\!l$ for every
$i\!\in\!\mathcal I_{k-1}$ is a \emph{candidate birth}. Under the
simplest birth model, each mask either corresponds to a newly-appearing
object (prior intensity $\lambda_b$, measured in units of expected
births per frame) or to clutter (intensity $\lambda_c$). Treat the
detector score $\hat s_{k,l}\!\in\![0,1]$ as a calibrated likelihood
ratio between the two: the unnormalised likelihood of the mask under
the ``new object'' hypothesis is proportional to $\hat s_{k,l}$, and
under the ``clutter'' hypothesis proportional to $1$ (independent of the
score). Bayes gives the posterior existence probability of the new
track,
\begin{equation}
   r^{(i_{\mathrm{new}})} \;=\;
   \frac{\lambda_b\,\hat s_{k,l}}
        {\lambda_b\,\hat s_{k,l} + \lambda_c},
   \label{eq:birth_r}
\end{equation}
which reduces to $\hat s_{k,l}$ when $\lambda_b=\lambda_c$ and smoothly
interpolates between a pure-score confidence ($\lambda_b\!\gg\!\lambda_c$)
and a clutter-dominated fallback ($\lambda_c\!\gg\!\lambda_b$). The
other parameters of the fresh Bernoulli are initialised from the
measurement itself:
\begin{equation}
   \ell^{(i_{\mathrm{new}})} = \hat\ell_{k,l},\quad
   T^{(i_{\mathrm{new}})}    = \hat T_{k,l},\quad
   P^{(i_{\mathrm{new}})}    = \hat R_{k,l}.
   \label{eq:birth_state}
\end{equation}
Since there is no pose prior on a freshly-born object, the posterior
covariance collapses to the measurement noise \cite{BarShalom2011}: this
is the Bayesian one-step update of an improper uniform prior by a
Gaussian likelihood. The new index $i_{\mathrm{new}}$ is added to
$\mathcal I_k$; $\lambda_b$ is a calibration constant per environment.

A track is \emph{reported} to downstream consumers (planning, shape
fusion, scene-graph computation) iff $r^{(i)}_{k|k}\ge r_{\text{conf}}$
(e.g.\ $0.5$); below this threshold the track is \emph{tentative} --
still predicted, associated, and updated, but masked from the report
$\widehat{\mathcal X}_k$.

%-----------------------------------------------------------------------------
\section{Visibility predicate $p_v$}\label{sec:vis}
%-----------------------------------------------------------------------------

We adopt a bounded, deterministic approximation of the true
observer-dependent visibility. Let $K\!\in\!\R^{3\times 3}$ be the
camera intrinsic matrix and
$\Omega_{\text{img}}\!=\![0,W]\!\times\![0,H]\!\subset\!\R^2$
the image domain. Define the camera projection
\begin{equation}
   \pi_K:\SE\to\R^2\cup\{\infty\},\qquad
   \pi_K(T) \;=\;
   \begin{cases}
      (Kp)_{1:2}/(Kp)_3, & p=T_{1:3,4},\ (Kp)_3>0,\\
      \infty,            & \text{otherwise},
   \end{cases}
\end{equation}
i.e.\ the pinhole projection of the object-origin in the camera frame
onto the image plane (with depth-positive gate). For tracks $i\!\ne\!j$,
let $A^{(i)}_k\!\subseteq\!\Omega_{\text{img}}$ be the image-space
axis-aligned bounding box of track $i$'s shape summary
$\mathcal G^{(i)}_{k-1}$ projected through
$T_{co}^{(i)}\!=\!T_{bc}^{-1}T^{(i)}_{k|k-1}$ (the base-frame mean
lifted to the camera frame via the fixed kinematic $T_{bc}$;
localization plays no role here), and $\bar d^{(i)}_k$ its mean
projected depth. Define the occlusion
coefficient
\begin{equation}
   \rho_{ij,k} \;\triangleq\;
   \mathbb 1\bigl[\bar d^{(j)}_k < \bar d^{(i)}_k\bigr]\,
   \frac{|A^{(j)}_k\cap A^{(i)}_k|}{|A^{(i)}_k|} \;\in\;[0,1],
   \label{eq:rho}
\end{equation}
which is the fraction of $i$'s projected footprint that is both covered
by $j$ in the image and \emph{in front of} $i$ along the viewing ray.
Distant or non-overlapping $j$ contribute $\rho_{ij,k}\!=\!0$. Then
\begin{equation}
   p_v^{(i)} \;=\;
   \mathbb 1\!\bigl[\pi_K\!\bigl(T_{bc}^{-1}\,T^{(i)}_{k|k-1}\bigr)\in\Omega_{\text{img}}\bigr]
   \;\cdot\;\prod_{j\in\mathcal I_{k-1}\setminus\{i\}}\bigl(1-\rho_{ij,k}\bigr).
   \label{eq:pv}
\end{equation}
$p_v^{(i)}\!=\!0$ when the projected centroid leaves the image, and the
product shrinks with each in-front occluder. \eqref{eq:pv} evaluates in
$O(|\mathcal I_{k-1}|^2)$ on bounding boxes; rendering-based refinement
(pixel-exact depth buffers) is a drop-in upgrade that does not change
any other equation.

%-----------------------------------------------------------------------------
\section{Pruning and output}\label{sec:emit}
%-----------------------------------------------------------------------------

After update, prune any track with $r^{(i)}_{k|k}<r_{\min}$
(e.g.\ $10^{-3}$). A track that survives a long sequence of
unobserved-but-expected frames sees its existence collapse via repeated
application of \eqref{eq:r_miss}, and is removed naturally without an
ad-hoc M-of-N counter.

The output at frame $k$ is the labelled MAP estimate -- in
\emph{base frame}, the native frame of the recursion:
\begin{equation}
   \widehat{\mathcal X}_k = \Bigl\{\,
       (i,\;\ell^{(i)},\;T^{(i)}_{k|k},\;P^{(i)}_{k|k},\;r^{(i)}_{k|k})\;:\;
       i\in\mathcal I_k,\ r^{(i)}_{k|k}\ge r_{\text{conf}}
   \,\Bigr\},
   \label{eq:emit_base}
\end{equation}
where $T^{(i)}_{k|k}=T_{bo,k|k}^{(i)}$ and
$P^{(i)}_{k|k}=P_{bo,k|k}^{(i)}$ are object-in-base.

\paragraph{World-frame composition (the only place $\Sigma_{wb,k}$ enters).}
Downstream consumers that work in world frame (planning, scene-graph
visualisation across episodes) compose with the current localization
posterior:
\begin{align}
   T_{wo,k}^{(i)} &= T_{wb,k}\,T_{bo,k|k}^{(i)},
       \label{eq:wo_mean}\\
   \Sigma_{wo,k}^{(i)} &= \operatorname{Ad}\!\bigl((T_{bo,k|k}^{(i)})^{-1}\bigr)\,
                          \Sigma_{wb,k}\,
                          \operatorname{Ad}\!\bigl((T_{bo,k|k}^{(i)})^{-1}\bigr)^{\!\top}
                       + P_{bo,k|k}^{(i)},
       \label{eq:wo_cov}
\end{align}
under the working assumption $\delta_{wb}\!\perp\!\delta_{bo}$ (a
reasonable approximation given base-frame storage decouples them
through the recursion). Eq.~\eqref{eq:wo_cov} lower-bounds $\Sigma_{wo}$
by the projected $\Sigma_{wb}$ -- exactly the desired physics for a
static object observed many times from the same base. World-frame
output is for reporting only; it is never re-ingested into the
recursion.

\bigskip
The single-frame composition above is the cheap default; it uses only
the current $T_{wb,k}$ and is blind to historical SLAM updates. A more
accurate world-frame estimate that absorbs the entire pose history
(loop closures included) is the chain smoother of \S\ref{sec:chain}.

%-----------------------------------------------------------------------------
\section{Per-track observation chain (canonical world-frame state)}
\label{sec:chain}
%-----------------------------------------------------------------------------

The Bernoulli-EKF of \S\ref{sec:ekf} is a one-step Markov filter:
$T^{(i)}_{k|k}$ depends on the previous posterior and the current
measurement only. That is the right tradeoff for the per-frame
association / existence loop, but it has a structural limitation: the
control input $u_k$ uses only the SLAM mean $T_{wb,k}$ (and its single
predecessor), so any retroactive change to the localization
trajectory -- a loop closure, a bundle-adjustment refinement, an
absolute-pose correction from a fiducial -- cannot be folded back
into the per-track posteriors that have already absorbed the old
$T_{wb}$ stream. The accumulated $P^{(i)}_{k|k}$ is wrong but
unrecoverable.

\paragraph{Chain definition.} For each track $i$ we additionally
maintain an append-only \emph{observation chain}
\begin{equation}
   \mathcal H^{(i)}_k \;=\;
   \Bigl\{\bigl(k_n,\;\hat T^{co,(i)}_{k_n,l},\;\hat R^{co,(i)}_{k_n,l}\bigr)
          \;:\; n=1,\dots,N^{(i)}_k\Bigr\},
   \label{eq:chain}
\end{equation}
where each entry records (a) the frame index $k_n$ at which track $i$
was matched to a detection, (b) the camera-frame ICP pose
$\hat T^{co,(i)}_{k_n,l}$ produced by Step~1 of \S\ref{subsec:bottleneck},
and (c) the camera-frame ICP covariance $\hat R^{co,(i)}_{k_n,l}$.
The chain is \emph{purely perception}: it does not store any
localization quantity. The fixed kinematic $T_{bc}$ promotes each
entry to base frame at zero variance (eq.~\ref{eq:meas_compose}).

\paragraph{Update rules.}
\begin{itemize}[leftmargin=2em,topsep=2pt,itemsep=0pt]
\item \textbf{Birth} (\S\ref{sec:birth}): initialise
      $\mathcal H^{(i_{\mathrm{new}})}_k = \{(k,\hat T^{co}_{k,l},\hat R^{co}_{k,l})\}$.
\item \textbf{Match} (\S\ref{sec:update}, matched branch): append
      $(k,\hat T^{co}_{k,l},\hat R^{co}_{k,l})$ to $\mathcal H^{(i)}_{k-1}$.
\item \textbf{Miss} (\S\ref{sec:update}, miss branch): chain unchanged.
\item \textbf{Prune} (\S\ref{sec:emit}): drop the chain.
\end{itemize}
Memory grows linearly with the matched-frame count of each track.
For long-running scenes a sliding window or a marginalize-into-prior
operation (replace the oldest $M$ entries with their fused Gaussian)
keeps the chain bounded; for typical tabletop episodes
$N^{(i)}_k\!\le\!10^3$ and no compression is needed.

\paragraph{Post-hoc world-frame composition.}
Given the current localization posterior over the pose history
$p\bigl(\{T_{wb,k}\}_{k=1:K}\bigr)$ -- its Gaussian posterior in the
filter case, the full pose-graph posterior when SLAM has loop
closures -- the world-frame pose of track $i$ is the smoothing
estimate
\begin{equation}
   \bigl(\hat T_{wo}^{(i)}, \hat\Sigma_{wo}^{(i)}\bigr)
   \;=\; \operatorname{Smooth}\!\bigl(\mathcal H^{(i)}_k,\,T_{bc},\,
            p(\{T_{wb,k}\}_{k=1:K})\bigr),
   \label{eq:chain_smooth}
\end{equation}
which we instantiate as a weighted Karcher-mean fusion of the per-entry
world-frame poses
\begin{align}
   \hat T_{wo,n}^{(i)} &\;=\; T_{wb,k_n}\,T_{bc}\,\hat T^{co,(i)}_{k_n,l},
       \label{eq:chain_lift_T}\\
   \hat R_{wo,n}^{(i)} &\;=\; \operatorname{Ad}\!\bigl(\hat T_{wo,n}^{(i)\,-1}\bigr)\,
                              \Sigma_{wb,k_n}\,
                              \operatorname{Ad}(\cdot)^{\!\top}
                            + \operatorname{Ad}(T_{bc})\,\hat R^{co,(i)}_{k_n,l}\,
                              \operatorname{Ad}(T_{bc})^{\!\top},
       \label{eq:chain_lift_R}
\end{align}
with weights $w_n\!\propto\!\Tr(\hat R_{wo,n}^{(i)})^{-1}$ in the simplest
information-weighted form. Iterated as a Gauss-Newton step on $\SE$
(re-linearise the innovations $\Log\!\bigl(\hat T_{wo,\mathrm{ref}}^{-1}
\hat T_{wo,n}\bigr)$, take the weighted mean correction, repeat until
convergence) this is a small per-track batch problem solved in
$O(N^{(i)}_k)$ time.

\paragraph{What this buys.}
\begin{enumerate}[leftmargin=2em,topsep=2pt,itemsep=0pt]
\item \textbf{Loop-closure consistency.} When SLAM revises the pose
      history (a bundle adjustment shifts $T_{wb,k_n}$ for some
      $k_n\!\in\!\mathcal H^{(i)}_k$), the chain is unchanged.
      Eq.~\eqref{eq:chain_smooth} immediately reflects the new pose
      history without re-running the EKF.
\item \textbf{Correct $\Sigma_{wb}$ correlation.} Two observations at
      the same base pose share the same $\Sigma_{wb,k}$ (they were
      taken from the same camera). The chain composition naturally
      preserves that correlation through the per-entry adjoint of
      eq.~\eqref{eq:chain_lift_R}; the EKF's one-step Markov form
      cannot.
\item \textbf{Compatibility with object-SLAM backends.} The chain is
      precisely the set of factors a downstream factor-graph
      smoother (GTSAM, iSAM2) needs: each entry is a binary factor
      between a base-pose node and an object-landmark node. Migration
      to a joint pose+object SLAM is a bookkeeping change, not an
      algorithmic one.
\end{enumerate}

\paragraph{Relation to the Bernoulli-EKF.} The two representations
co-exist. Per-frame association, Bernoulli existence updates, and
birth / prune all read $(T^{(i)}_{k|k}, P^{(i)}_{k|k})$ from the
EKF (the \emph{fast tier}); world-frame consumers and SLAM-aware
queries read $(\hat T_{wo}^{(i)}, \hat\Sigma_{wo}^{(i)})$ from the
chain (the \emph{canonical world-frame state}). When a SLAM update
lands, the chain composition is recomputed and the result can
optionally be projected back to base frame
($T_{bo}^{(i)} = T_{wb,k}^{-1}\hat T_{wo}^{(i)}$) and \emph{injected}
into the EKF as a re-initialisation; this is the periodic
reconciliation pattern already used by the slow tier in the
existing implementation.

%-----------------------------------------------------------------------------
\section{Frame-by-frame algorithm}
%-----------------------------------------------------------------------------

\begin{enumerate}[leftmargin=2em]
   \item \textbf{Predict.} For each $i\in\mathcal I_{k-1}$: pick $u^{(i)}_k$ per eq.~\eqref{eq:control} ($\Delta T_{bg,k}$ for manipulation set, $\Delta T_{wb,k}^{-1}$ otherwise); apply state predict via \eqref{eq:ekf_T_pred}--\eqref{eq:ekf_P_pred}; advance existence via \eqref{eq:bern_pred_r}.
   \item \textbf{Measure.} For every mask $m_{k,l}$ back-project to a \emph{camera-frame} point cloud $\mathcal C_{k,l}$. For each candidate track $i$ passing a coarse gate, run colored ICP against $\mathcal G^{(i)}_{k-1}$ to obtain $(\hat T^{co}_{k,l},\hat R^{co}_{k,l})$, then promote to base via \eqref{eq:meas_compose}: $\hat T_{k,l}=T_{bc}\hat T^{co}_{k,l}$, $\hat R_{k,l}=\operatorname{Ad}(T_{bc})\hat R^{co}_{k,l}\operatorname{Ad}(T_{bc})^{\!\top}$. Compute $d^2_{il}$ via \eqref{eq:innov_pair} and $C_{il}$ via \eqref{eq:cost}.
   \item \textbf{Associate.} Solve \eqref{eq:hung} for $\pi^{\!*}$.
   \item \textbf{Update each track $i$.}
         \begin{itemize}[topsep=0pt]
            \item if $\pi^{\!*}(i)=l$ and $w_{il}\!>\!0$: apply \eqref{eq:r_assoc} and EKF \eqref{eq:ekf_T_upd}--\eqref{eq:ekf_P_upd} with $R_{il}/w_{il}$; set $\tau^{(i)}_k\!\leftarrow\!\tau_{k,l}$; \emph{append} $(k,\hat T^{co}_{k,l},\hat R^{co}_{k,l})$ to the observation chain $\mathcal H^{(i)}_{k-1}$ (\S\ref{sec:chain});
            \item else: compute $p_v^{(i)}$ via \eqref{eq:pv} and apply \eqref{eq:r_miss} (leaving $\tau^{(i)}_k\!=\!\tau^{(i)}_{k-1}$ and $\mathcal H^{(i)}$ unchanged).
         \end{itemize}
   \item \textbf{Birth.} For each unassigned mask: centroid-initialised $\hat T$ and inflated $\hat R$, then \eqref{eq:birth_r}, \eqref{eq:birth_state}, $\tau^{(i_{\mathrm{new}})}_k\!\leftarrow\!\tau_{k,l}$, and $\mathcal H^{(i_{\mathrm{new}})}_k\!=\!\{(k,\hat T^{co}_{k,l},\hat R^{co}_{k,l})\}$.
   \item \textbf{Prune.} Drop $i$ with $r^{(i)}_{k|k}<r_{\min}$.
   \item \textbf{Emit.} Report base-frame $\widehat{\mathcal X}_k$ via \eqref{eq:emit_base}; world-frame consumers compose with $T_{wb,k},\Sigma_{wb,k}$ via \eqref{eq:wo_mean}--\eqref{eq:wo_cov}.
\end{enumerate}

%-----------------------------------------------------------------------------
\section{Reduction to existing code}
%-----------------------------------------------------------------------------

\begin{itemize}[leftmargin=2em]
   \item Bernoulli state $(r^{(i)},\ell^{(i)},T^{(i)},P^{(i)},\tau^{(i)})$
         replaces the current per-track entry; two extra scalar fields
         added to \texttt{rbpf\_state.RBPFState}: existence $r$ and
         last-matched SAM2 tracklet ID $\tau$.
   \item Predict \eqref{eq:bern_pred_r}, \eqref{eq:ekf_T_pred}--\eqref{eq:ekf_P_pred} edits \texttt{rbpf\_state.predict\_objects};
         the saturation projector $\Phi$ replaces the bucket-switching
         in \texttt{process\_noise\_for\_phase}.
   \item Cost matrix \eqref{eq:cost}, assignment \eqref{eq:hung},
         and per-pair update \eqref{eq:r_assoc} + Huber-weighted EKF
         update replace the hard upstream-ID lookup in
         \texttt{orchestrator.\_fast\_tier} with a soft, ID-bonused
         Mahalanobis assignment: SAM2 IDs still drive identity when
         consistent, but association degrades gracefully to pure
         geometry when SAM2 drops an ID.
   \item \eqref{eq:r_miss} replaces the integer \texttt{miss\_streak};
         pruning at $r_{\min}$ replaces the M-of-N deletion test.
   \item $p_v^{(i)}$ via \eqref{eq:pv} adds a new \texttt{visibility.py}
         module ($\sim$60 LOC, depends only on the fixed $T_{bc}$ and
         the per-object shape summaries -- no SLAM input).
   \item Base-frame storage of $T^{(i)}, P^{(i)}$ aligns with the
         existing \texttt{pose\_update.gaussian\_state.GaussianState}
         convention; world-frame composition lives in
         \texttt{collapsed\_object\_world}. The RBPF Bernoulli path is
         being migrated to base-frame storage to match.
\end{itemize}

\paragraph{Calibrated parameters.} The recursion is parameterised by
seven scalars,
\[
   p_s\!\in\!(0,1],\quad
   p_d\!\in\!(0,1],\quad
   \alpha\!\ge\!0,\quad
   \lambda_c\!>\!0,\quad
   \lambda_b\!>\!0,\quad
   r_{\text{conf}}\!\in\!(0,1),\quad
   r_{\min}\!\in\!(0,r_{\text{conf}}),
\]
one symmetric positive-definite matrix $P_{\max}\!\in\!\R^{6\times6}$
(the covariance saturation cap of \eqref{eq:phi}), and two
$\chi^2_6$-quantiles
$G_{\text{in}}\!=\!\chi^2_6(0.95)\!\approx\!12.59$,
$G_{\text{out}}\!=\!\chi^2_6(0.9997)\!\approx\!25$. Defaults for static
manipulation scenes: $p_s\!=\!1$ (rigid objects do not spontaneously
disappear), $p_d\!=\!0.9$, $\alpha\!=\!4.4$ (mild tie-break toward
SAM2 ID continuity; calibrate on a held-out sequence), $\lambda_c$
estimated from a few static frames as the average per-frame detection
count divided by image area, $\lambda_b\!=\!\lambda_c$,
$r_{\text{conf}}\!=\!0.5$, $r_{\min}\!=\!10^{-3}$,
$P_{\max}\!=\!\operatorname{diag}(\sigma^2_{\text{trans}}\!\cdot\!\bm 1_3,\,\sigma^2_{\text{rot}}\!\cdot\!\bm 1_3)$
with $\sigma_{\text{trans}}\!=\!0.25\,\mathrm{m},\,\sigma_{\text{rot}}\!=\!\pi/4$.
Setting $p_s\!<\!1$ models a known per-frame physical-removal hazard
(e.g.\ unstable platforms, transient occluding clutter) and is purely
domain-specific; setting $\alpha\!=\!0$ switches off the SAM2 ID prior
entirely and recovers the pure-geometry cost \eqref{eq:cost}.

%-----------------------------------------------------------------------------
\section{Approximations relative to PMBM}\label{sec:approx}
%-----------------------------------------------------------------------------

The full PMBM posterior of \cite{GarciaFernandez2018} is
\[
   f_k(\mathcal T)
   = \!\!\!\!\sum_{Y\uplus W=\mathcal T}\!\!\!\!
       \underbrace{f^p(Y)}_{\text{PPP, undetected}}\;
       \underbrace{f^{mbm}(W)}_{\text{MB mixture, detected}},
   \quad
   f^{mbm}(W)\propto \sum_j\sum_{X_1\uplus\dots\uplus X_n=W}\,\prod_i w_{j,i}\,f_{j,i}(X_i).
\]
Our recursion is its small-cardinality, low-clutter, single-hypothesis
limit. We list below every place where it deviates and label each
deviation as either an \emph{approximation} (lossy collapse of an
exact PMBM construct), a \emph{modelling choice} (different but
internally Bayesian model), or a \emph{linearisation} (first-order
truncation).

\begin{enumerate}[leftmargin=2em,label=\textbf{A.\arabic*},itemsep=2pt]

\item \textbf{Single-hypothesis collapse via Hungarian assignment}
      \emph{(approximation)}. We retain only the MAP global hypothesis
      $j^{\!*}$ from $\sum_j w_j(\cdot)$, found by linear assignment
      \eqref{eq:hung}. This is the JIPDA reduction of PMBM
      (\cite{GarciaFernandez2018} Sec.\,IV-B): conditional on
      $j^{\!*}$ being the true assignment, all subsequent updates are
      exact Bayesian. Equivalent to PMBM's mixture only when the MAP
      assignment is unambiguous.

\item \textbf{Marginalised single-MB existence form}
      \emph{(equivalence at marginal level)}. PMBM sets $r^+\!=\!1$ on
      a successful association and pushes the ambiguity into the
      hypothesis weight $w_{j,i}$. After collapse to the MAP, those
      separately-tracked quantities marginalise to our
      \eqref{eq:r_assoc}: $r^+ = p_d\mathcal L\,r/(p_d\mathcal L\,r + \lambda_c(1-r))$.
      The misdetection update \eqref{eq:r_miss} is algebraically
      identical to PMBM's eq.\,(31) (with state-dependent
      $\tilde p_d^{(i)}$).

\item \textbf{No PPP for never-detected targets} \emph{(modelling
      choice)}. We do not maintain the Poisson intensity $\mu(\cdot)$
      that PMBM uses to model targets that exist but have not yet
      produced a measurement. Equivalent to assuming every object in
      the scene eventually produces at least one detection, which holds
      for our manipulation domain (objects are deliberately placed by a
      human or robot and not invisible at first appearance). Loses
      PMBM's spatial prior over unseen states.

\item \textbf{Categorical birth from detection score}
      \emph{(modelling choice)}. We replace PMBM's birth integral
      $e(z)\!=\!p_d\!\int p(z|y)\mu(y)dy$ with a per-mask Bayes
      decision over a binary alternative ``new object vs.\ clutter,''
      using the detector score $\hat s_{k,l}$ as the likelihood ratio
      and rate-constants $\lambda_b,\lambda_c$ as the prior odds:
      \eqref{eq:birth_r}. Internally Bayesian under this model; not
      spatially adaptive.

\item \textbf{Multi-Bernoulli factorisation across indices}
      \emph{(modelling choice)}. We assume the per-index Bernoullis
      are statistically independent. PMBM also makes this assumption
      \emph{within} each global hypothesis $j$; we make it
      \emph{globally} after the MAP collapse.

\item \textbf{Constant survival probability with $p_s\!=\!1$ default}
      \emph{(modelling choice)}. We follow standard PMBM in keeping
      $p_s$ observer-independent; for static manipulation scenes the
      natural value is $p_s\!=\!1$, removing predict-step decay
      entirely. Visibility-dependent ``apparent'' survival is recovered
      automatically through the state-dependent $\tilde p_d^{(i)}$ in
      the missed update, which is admissible PMBM.

\item \textbf{EKF on $\SE$ with $H\!=\!I_6$} \emph{(linearisation)}.
      Process and measurement noise are right-trivialised in the body
      tangent; the state composition and the innovation use BCH
      truncated at first order. Exact at $\nu_k\!=\!\bm 0$; smoothly
      degrading with $\|\nu_k\|$. Standard EKF-on-Lie-group practice.

\item \textbf{Covariance saturation $\Phi$} \emph{(modelling choice)}.
      The cap on $\Tr P$ at $\Tr P_{\max}$ in \eqref{eq:phi} is a
      pragmatic stabiliser for arbitrarily long missing-observation
      streaks; it equals an additive Gaussian prior $P_{\max}$ that
      is not in standard PMBM but is consistent with a maximum-entropy
      regularisation.

\end{enumerate}

\paragraph{What this buys, what this costs.}
Approximations A.1 and A.2 are exactly the JIPDA reduction of PMBM,
known to be a tight approximation when MAP associations are
unambiguous and lossy when crossing tracks induce hypothesis ambiguity.
Modelling choices A.3 -- A.6 specialise the standard PMBM model to
our domain and are reversible: restoring the PPP, Murty K-best
mixture, spatially-adaptive birth, or per-class $p_s$ would re-recover
full PMBM at proportional computational cost. Linearisations A.7 --
A.8 are inherent to the EKF-on-Lie-group and covariance-bounded
layer; they do not change the multi-target structure of the
recursion.

\begin{thebibliography}{9}
\bibitem{Williams2015}
J.~L. Williams. ``Marginal multi-Bernoulli filters: RFS derivation of
MHT, JIPDA, and association-based MeMBer.''
\textit{IEEE Transactions on Aerospace and Electronic Systems},
51(3): 1664--1687, 2015.

\bibitem{GarciaFernandez2018}
\'A.~F. Garc\'{\i}a-Fern\'{a}ndez, J.~L. Williams, K. Granstr\"{o}m,
L. Svensson. ``Poisson multi-Bernoulli mixture filter: Direct derivation
and implementation.'' \textit{IEEE Transactions on Aerospace and
Electronic Systems}, 54(4): 1883--1901, 2018.

\bibitem{BarShalom2011}
Y. Bar-Shalom, P.~K. Willett, X. Tian.
\textit{Tracking and Data Association}. YBS Publishing, 2011.

\bibitem{MurrayLiSastry1994}
R.~M. Murray, Z. Li, S.~S. Sastry. \textit{A Mathematical Introduction
to Robotic Manipulation}. CRC Press, 1994.
\end{thebibliography}

\end{document}
